\subsection{可测性判别法}

当 F 是拓扑向量空间时，每个取值于 F 的、以可测集为底的阶梯函数（§4, No.~9, 定义 4）显然是可测的（No.~3, 定理 1 的推论 3）；这样的函数 f 只取有限多个值，且对每个 \(y \in F\)，\(f(y)\) 是可测的。更一般地，设 F 是任意拓扑空间，f 是 X 到 F 中只取有限多个不同值 \(a_i\)（ \(1 \leqslant i \leqslant m\) ）的一个映射；若诸集合 \(A_i = f(a_i)\) 可测，则函数 f 可测。因为，对每个紧集 K 以及诸集合 \(A_i \cap K\) 中的每一个，存在可忽略集 \(N_i \subset A_i \cap K\) 以及 \(A_i \cap K \cap C N_i\) 的由紧集序列（ \(K_{in}\) ）构成的一个分割；由于 K 是诸集合 \(A_i \cap K\) 的并，且 f 到每个 \(K_{in}\) 上的限制是常数、从而是连续的，故 f 可测。我们约定，按措辞的一种滥用，称 X 到 F 中的映射 f 为可测阶梯函数，如果它只取有限多个值，且对每个 \(y \in F\)，\(f(y)\) 是可测的。

定理 3. --- 要使 X 到可度量化空间 F 中的映射 f 可测，必须且只需：对每个紧集 \(K \subset X\)，存在取值于 F 的可测阶梯函数的一个序列 \((g_{n})\)，使得 \(g_{n}(x)\) 趋于 \(f(x)\)（对几乎所有 \(x \in K\)）。

条件是充分的，由 Egoroff 定理与局部化原理即得。我们来证明它是必要的：由假设，存在可忽略集 \(N \subset K\) 以及由紧集组成的 K - N 的一个分割 \((K_{m})\)，使得 f 到每个 \(K_{m}\) 上的限制是连续的。为定义序列 \((g_{n})\)，只需按如下方式进行：设 d 是与 F 的拓扑相容的一个距离；对每个 \(K_{i}\)（指标满足 \(i \leqslant n\) 者），存在 \(K_{i}\) 分成可积集 \(A_{ij}\)（ \(1 \leqslant j \leqslant q_{i}\) ）的一个有限分割

使得 \(f\) 在每个 \(A_{ij}\) 上的振幅 \(\leqslant 1/n\)（§4, No.~9, 引理）；取 \(g_n\) 在每个 \(A_{ij}\) 上为常数，等于 \(f\) 在该集合中的一个值（对 \(1 \leqslant i \leqslant n\)、\(1 \leqslant j \leqslant q_i\)），等于一个固定元素 \(a \in F\)（对 \(X\) 中不属于任何 \(A_{ij}\) 的每个点）。显然，序列 \((g_n(x))\) 收敛于 \(f(x)\)（在 \(K\) 中每个不属于 \(N\) 的点上）。

推论 1. --- 设 f 是 X 到 Banach 空间 F 中的一个可测映射；对每个紧集 K ⊂ X，存在可测阶梯函数的一个序列 (gₙ)，其支集都含于 K，使得对所有 x ∈ X 有 \textbar gₙ(x)\textbar{} ≤ \textbar f(x)\textbar{}，且对几乎所有 x ∈ K，gₙ(x) 趋于 f(x)。

沿用定理 3 证明中的记号，并用 \(\mathbf{a}_{ij}\) 表示 \(\mathbf{f}\) 在 \(\mathrm{A}_{ij}\) 中的一个值，只需取 \(\mathbf{g}_n\) 在 \(\mathrm{A}_{ij}\) 中的值为：当 \(|\mathbf{a}_{ij}| \leqslant 1/n\) 时取点 0，否则取点 \(\mathbf{a}_{ij}(1 - 1/(n|\mathbf{a}_{ij}|))\)；最后，令 \(\mathbf{g}_n(x) = 0\) 于诸 \(\mathrm{A}_{ij}\)（ \(1 \leqslant i \leqslant n\)，\(1 \leqslant j \leqslant q_i\) ）的并的余集上。

推论 2. --- 设 X 是一个无穷远处可数的局部紧空间。若 f 是 X 到可度量化空间 F 中的可测映射，则存在取值于 F 的可测阶梯函数的一个序列 \((g_{n})\)，使得 \(g_{n}(x)\) 趋于 \(f(x)\)（对几乎所有 \(x \in X\)）。

若 X 是紧集的递增序列 \((\mathbf{A}_{n})\) 的并，则诸非空的集合 \(A_{n}-A_{n-1}\) 构成 X 分成可积集的一个分割；于是，存在可忽略集 \(N\subset X\) 以及由紧集序列 \((\mathrm{K}_{n})\) 构成的 CN 的一个分割，使得 f 到每个 \(K_{n}\) 上的限制是连续的；此后证明可毫无改动地像定理 3 那样完成。

命题 7. --- 设 f 是 X 到拓扑空间 F 中的一个可测映射；F 中每个闭（相应地，开）集在 f 下的逆像是可测的。

只需对 F 中闭集 A 的逆像 \(f(A)\) 进行证明。设 K 是 X 的一个紧子集；存在可忽略集 \(N \subset K\) 以及由紧集组成的 K - N 的一个分割 \((K_n)\)，使得 f 到每个 \(K_n\) 上的限制是连续的。于是交 \(K_n \cap f(A)\) 是闭集 A 在 f 到 \(K_n\) 上的限制下的逆像；因此它是 \(K_n\) 中的闭集，从而是紧的。于是集合 \(K \cap f(A)\) 是可忽略集 \(N \cap f(A)\) 与诸紧集 \(K_n \cap f(A)\) 的并，这就证明了 \(f(A)\) 是可测的。

定理 4. --- 设 F 是一个可度量化空间，d 是 F 上与该拓扑相容的一个距离。要使 X 到 F 中的映射 f 可测，必须且只需它满足下列两个条件：

\begin{enumerate}
\def\labelenumi{\alph{enumi})}
\item
  在 \(f\) 下，\(F\) 的每个闭球的逆像是可测的；
\item
  对每个紧集 \(\mathrm{K} \subset \mathrm{X}\)，存在一个可数子集 \(\mathrm{H}\)（\(\mathrm{F}\) 的），使得 \(f(x) \in \overline{\mathrm{H}}\) 对几乎所有 \(x \in \mathrm{K}\) 成立。
\end{enumerate}

条件 a) 是必要的，由命题 7 即得；另一方面，在定理 3 的记号下，取 H 为由所有函数 \(g_{n}\) 的值组成的可数集，条件 b) 即得到满足。

现在来证明条件 a) 与 b) 是充分的。设 K 是 X 的任意紧子集；存在 K 的一个可忽略子集 N，使得 \(f(\mathrm{K}-\mathrm{N})\) 含于 F 的一个可数点集的闭包中，我们将该点集排成一个序列 \((a_{n})\)。设 \(A_{n,p}\) 是点 \(x \in K-N\)（使得 \(d(f(x), a_{n}) \leqslant 1/p\) 者）所成之集。由条件 a) 可知 \(A_{n,p}\) 是可测的。对固定的 p，递归地定义集合序列 \(B_{n,p} \subset K-N\) 如下：

\[
\mathrm{B} _ {1, p} = \mathrm{A} _ {1, p} \quad \text {and} \quad \mathrm{B} _ {n + 1, p} = \mathrm{A} _ {n + 1, p} \cap \mathbb {C} \Big (\bigcup_ {k \leqslant n} \mathrm{A} _ {k, p} \Big);
\]

诸集合 \(B_{n,p}\) 是可测的，且其中非空者构成 K-N 的一个分割。设 \(g_{m,p}\) 是这样的函数：等于 \(a_{i}\) 于集合 \(B_{i,p}\) 上（对 \(1 \leqslant i \leqslant m\)），在这些集合的并的余集上等于一个常数 \(b \in F\)；\(g_{m,p}\) 是一个可测阶梯函数；当 m 趋于无穷时，\(g_{m,p}\) 逐点收敛于函数 \(f_{p}\)，后者等于 \(a_{n}\) 于 \(B_{n,p}\)（ \(n \geqslant 1\) ）上，在 N∪CK 上等于 b，因此（定理 2）\(f_{p}\) 是可测的。当 p 趋于无穷时，\(f_{p}(x)\) 趋于 \(f(x)\)（对每个 \(x \in K-N\)），而对 \(x \in N \cup CK\) 趋于 b；因此诸 \(f_{p}\) 的极限是可测的，再由局部化原理即知 f 本身是可测的。

注. --- 1) 仅条件 a) 不足以使 \(f\) 可测（习题 7）。

因为，若确实如此，则对每个 \(b \in \overline{R}\)，使 \(f(x) \geqslant b\) 的 x 所成之集是可测的，因为它是使 \(f(x) \geqslant a\) 的 x 所成之集（它们构成一个可数族）的交，其中 a 遍历 D 中 \(\leqslant b\) 的点所成之集。使 \(f(x) < b\) 的 x 所成之集是可测的，因为它是一个可测集的余集。其次，若 b 有限，则使 \(f(x) \leqslant b\) 的 x 所成之集是可测的，因为它是使 \(f(x) < b + 1/n\) 的 x 所成之集的交；而 \(\overline{f}(-\infty)\) 是可测的，因为当 n 遍历 Z 时，它是使 \(f(x) < n\) 的 x 所成之集的交。最后，f 下 \(\overline{R}\) 的每个闭区间的逆像是可测的，因为它是两个可测集的交，于是可以应用定理 4。

同样可以证明，只需对每个 \(a \in D\)，\(x\)（使得 \(f(x) > a\) 者）所成之集是可测的即可。

推论. --- 每个下半（相应地，上半）连续函数都是可测的。

因为若 \(f\) 是下半连续的，则 \(x \in \mathbf{X}\)（使得 \(f(x) \leqslant a\) 者）所成之集对每个 \(a \in \overline{\mathbf{R}}\) 都是闭的。

当 F 可度量化且紧时，命题 7 使我们能将定理 3 加强如下：

命题 9. --- 若 F 是一个可度量化的紧空间，则每个 X 到 F 中的可测映射 f 都是可测阶梯函数序列的（在整个 X 上的）一致极限。

设 \(d\) 是与 \(\mathbf{F}\) 的拓扑相容的一个距离。对每个正整数 \(n\)，存在有限多个点 \(a_{k} \in \mathbf{F}\)，使得闭球 \(\mathbf{B}_{k}\)（以 \(a_{k}\) 为中心、\(1/n\) 为半径）构成 \(\mathbf{F}\) 的一个覆盖；于是诸集合 \(\mathbf{A}_{k} = f(\mathbf{B}_{k})^{-1}\) 是可测的（命题 7）且构成 \(\mathbf{X}\) 的一个覆盖。因此（§4, No.~9, 引理）存在一个分割 \((\mathbf{C}_{i})\)（把 \(\mathbf{X}\) 分成有限多个可测集的），使得每个 \(\mathbf{A}_{k}\) 都是某些 \(\mathbf{C}_{i}\) 的并。设 \(c_{i}\) 是 \(\mathbf{C}_{i}\) 的一个点，并设 \(g_{n}\) 是等于 \(f(c_{i})\) 于 \(\mathbf{C}_{i}\) 上（对每个指标 \(i\)）的可测阶梯函数。显然，\(d(f(x), g_{n}(x)) \leqslant 2/n\) 对所有 \(x \in \mathbf{X}\) 成立。

命题 10. --- 设 F 是一个可分 Banach 空间，F' 是它的对偶，而 (a'\_n) 是 F' 的单位球中的一个弱稠密序列（TVS, III, §3, No.~4, 命题 6 的推论 2）。¹ 要使 X 到 F 中的映射 f 可测，必须且只需对每个 n，标量函数 x ↦ ⟨f(x), a'\_n⟩ 都是可测的。

条件显然是必要的（No.~3, 定理 1），我们来证明它是充分的；只需验证定理 4 的条件 a)，为此，由局部化原理，只需证明：对 X 的每个紧子集 K 以及 F 中每个以 a 为中心、r 为半径的闭球 \(B \subset F\)，集合 \(A = K \cap \mathbf{f}^{-1}(B)\) 都是可测的；而现在，对每个 \(z \in F\)，

\[
| \mathbf {z} | = \sup _ {n} | \langle \mathbf {z}, \mathbf {a} _ {n} ^ {\prime} \rangle | / | \mathbf {a} _ {n} ^ {\prime} |;
\]

于是 A 是 \(K\) 与由下式定义的诸集合的交

\[
\left| \left\langle \mathbf {f} (x), \mathbf {a} _ {n} ^ {\prime} \right\rangle - \left\langle \mathbf {a}, \mathbf {a} _ {n} ^ {\prime} \right\rangle \right| \leqslant r \left| \mathbf {a} _ {n} ^ {\prime} \right|;
\]

由于这些集合依假设是可测的，故 A 是可测的。

推论 1. --- 设 F 是一个 Banach 空间。要使 X 到 F 中的映射 f 可测，必须且只需它满足下列两个条件：

\begin{enumerate}
\def\labelenumi{\alph{enumi})}
\tightlist
\item
  对每个 \(\mathbf{a}' \in \mathrm{F}'\)，标量函数 \(x \mapsto \langle \mathbf{f}(x), \mathbf{a}' \rangle\) 是可测的；
\item
  对每个紧集 \(\mathrm{K} \subset \mathrm{X}\)，存在一个可数子集 \(\mathrm{H}\)（\(\mathrm{F}\) 的），使得 \(\mathbf{f}(x) \in \overline{\mathrm{H}}\) 对几乎所有 \(x \in \mathrm{K}\) 成立。
\end{enumerate}

这些条件的必要性由 No.~3 的定理 1 以及定理 4 得出。为证明条件是充分的，仍只需验证定理 4 的条件 a)。沿用命题 10 证明中的记号，我们可以（借助于 b)）在必要时于一个可忽略集上修改 f 之后，假设 \(\mathbf{f}(\mathrm{K}) \subset \overline{\mathrm{H}}\)，其中 H 是 F 的一个可数子集。若 V 是由集合 \(H \cup \{a\}\) 生成的 F 的闭线性子空间，则 V 是一个可分 Banach 空间，且 V 上每个连续线性形式都是某个形式 \(a' \in F'\) 的限制；于是与命题 10 中同样的推理表明 \(K \cap \overline{\mathbf{f}}(B)\) 是可测的。

推论 2. --- 设 F 是一个可度量化且可分的局部凸空间，并设 \(F'\) 是它的对偶。要使 X 到 F 中的映射 f 可测，必须且只需对每个 \(a' \in F'\)，标量函数 \(x \mapsto \langle \mathbf{f}(x), \mathbf{a}' \rangle\) 都是可测的。

我们可以把 F 视为 Banach 空间的一个可数乘积 \(\prod_{n}E_{n}\) 的子空间（TVS, II, §4, No.~3），并且可以假设 \(\operatorname{pr}_{n}(\mathbf{F})\) 在 \(E_{n}\) 中稠密，后者因而是可分的。这样，对每个 n，映射 \(pr_{n}\circ f\) 依据命题 10 是可测的，因此 f 依据 No.~3 的定理 1 是可测的。

命题 11. --- 设 F 是一个局部凸空间，它是可分且可度量的局部凸空间的一个序列 \(F_{n}\) 的直接极限，且 F 是诸 \(F_{n}\) 的并。设 \(F'\) 是 F 的对偶，赋以弱拓扑 \(\sigma(\mathrm{F}',\mathrm{F})\)。要使 X 到 \(F'\) 中的映射 f 可测，必须且只需：对每个 \(a \in F\)，标量函数 \(x \mapsto \langle \mathbf{a}, \mathbf{f}(x) \rangle\) 都是可测的。

条件是必要的，由 No.~3 的定理 1 即得；我们来证明它是充分的。首先假设 F 可度量化且可分，并设 D 是 F 中的一个可数稠密集。设 \((\mathrm{V}_{n})\) 是 F 中 0 的均衡凸开邻域的一个递减基本序列；极集 \(V_{n}^{\circ}\) 是等度连续的，且它们的并就是整个 \(F'\)。设 \(X_{n} = \overline{\mathbf{f}}(V_{n}^{\circ})\)；序列 \((\mathrm{X}_{n})\) 是递增的，且 \(X = \bigcup_{n} X_{n}\)；我们来证明每个 \(X_{n}\) 都是 \(\mu\)-可测的。事实上，\(D \cap V_{n}\) 在 \(V_{n}\) 中稠密；对每个 \(y \in D \cap V_{n}\)，设 \(S_{y}\) 是 \(x \in X\)（使得 \(|\langle y, f(x) \rangle| \leqslant 1\) 者）所成之集；假设蕴含每个 \(S_{y}\) 都可测，而 \(X_{n}\) 是诸 \(S_{y}\)（对 \(y \in D \cap V_{n}\)）组成的可数族的交。在此基础上，对 X 的每个紧子集 K 和每个 \(\varepsilon > 0\)，存在整数 n，使得 \(|\mu| (K - (K \cap X_{n})) \leqslant \varepsilon/4\)，然后存在紧子集 \(K_{1}\)（\(K \cap X_{n}\) 的），使得 \(|\mu| ((K \cap X_{n}) - K_{1}) \leqslant \varepsilon/4\)；最后，存在紧子集 \(K_{2}\)（\(K_{1}\) 的），使得 \(|\mu| (K_{1} - K_{2}) \leqslant \varepsilon/2\)，且在 \(K_{2}\) 上所有函数 \(\langle y, f \rangle\)（其中 \(y \in D\)）的限制都是连续的（No.~1, 命题 2）。由于集合 \(f(K_{2}) \subset f(X_{n}) \subset V_{n}^{\circ}\) 是等度连续的，\(\sigma(F', F)\) 在 \(f(K_{2})\) 上诱导的拓扑与在 D 上逐点收敛的拓扑相同（GT, X, §2, No.~4, 定理 1）；因此 f 到 \(K_{2}\) 上的限制是连续的，由此得到第一种情形下的断言。

现在过渡到一般情形。若 \(\mathbf{z}'\) 是 \(\mathbf{F}\) 上的一个连续线性形式，则其限制 \(\mathbf{z}_n'\)（到 \(\mathbf{F}_n\) 上的）是连续的；由于 \(\mathbf{F} = \bigcup_{n} \mathbf{F}_n\)，对偶 \(\mathbf{F}'\)（\(\mathbf{F}\) 的）可以（在代数意义上）等同于乘积 \(\prod_{n} \mathbf{F}_n'\) 的一个线性子空间，且此时 \(\operatorname{pr}_n(\mathbf{z}') = \mathbf{z}_n'\)。此外，由于 \(\mathbf{F}\) 的每个有限子集都含于某个 \(\mathbf{F}_n\)，拓扑 \(\sigma(\mathbf{F}', \mathbf{F})\) 恰是诸拓扑 \(\sigma(\mathbf{F}_n', \mathbf{F}_n)\) 的乘积拓扑所诱导的拓扑。在此基础上，若 \(\langle \mathbf{a}, \mathbf{f} \rangle\) 对每个 \(\mathbf{a} \in \mathbf{F}\) 可测，则 \(\langle \mathbf{a}_n, \operatorname{pr}_n \circ \mathbf{f} \rangle\) 对每个 \(n\) 与每个 \(\mathbf{a}_n \in \mathbf{F}_n\) 也可测，因为 \(\langle \mathbf{a}_n, \operatorname{pr}_n \circ \mathbf{f} \rangle = \langle \mathbf{a}_n, \mathbf{f} \rangle\)；于是证明的第一部分表明，\(\operatorname{pr}_n \circ \mathbf{f}\) 对每个 \(n\) 可测，因此 \(\mathbf{f}\) 也可测（No.~3, 定理 1）。

